\subsection*{An actual nonnilpotent finite quotient of the residual group}

\statement{B.62}{Proved}{The mod-three projective homology quotient kills the residual centralizer}
Let $P=P\Gamma_{0,12}$, and let
\[
 q:\Gamma_{0,[12]}\longrightarrow
       \operatorname{PSp}_{10}(\mathbb F_3)
                                                               \tag{B.62.1}
\]
be the projective homology action on the genus-five double
cover of the sphere branched over the twelve named points.
Its restriction to $P$ is surjective.  In the full profinite
completion of $P$,
\[
                \widehat q(\mathcal K)=1.                    \tag{B.62.2}
\]
Consequently, if $L_1,L_2\leq_f P$ both contain
$N_3=\ker(q|_P)$, then
\[
                 (H_T\cap L_1)(H_5\cap L_2)                 \tag{B.62.3}
\]
is closed in the full profinite topology of $P$.
The kernel $N_3$ is invariant under the full sphere
mapping-class group, including the given $A_5$.
The quotient controlled here is nonnilpotent; the assertion
does not identify $\mathcal K$, control all finite quotients,
or assert vanishing at higher $3$-power levels.

\begin{proof}
\emph{The geometric representation and its integral model.}
Send every peripheral loop of the twelve-punctured sphere
to the nonidentity element of $C_2$.
Because twelve is even, this defines a connected double
cover, which extends across the punctures to a branched
cover $\Sigma_5\to S^2$.
Every orientation-preserving mapping class permutes the
peripheral conjugacy classes and thus lifts to this cover.
The two lifts differ by its hyperelliptic involution, whose
action on integral first homology is $-I$.
Lifting isotopies shows that the resulting projective
symplectic action depends only on the sphere mapping
class.  This defines (B.62.1) after reduction modulo three.

With label $2$ at infinity, the disk cover has genus five
and one boundary component.  Capping its boundary preserves
first homology.  This representation is the nontrivial
summand of unreduced Burau at $t=-1$.
The branched-cover chain calculation is given by
\citet[Sections 2.2--2.3, Proposition 2.1]{DE_GambaudoGhys2005};
the odd-strand splitting and the action $-I$ of the full
twist are also explicitly recorded by
\citet[Introduction, ``Hyperelliptic Torelli vs Burau'']
{DE_BrendleMargalitPutman2015}.

Here is a fully specified integral convention for the
calculation.  The unreduced positive generator $\sigma_j$
has block
\[
              \begin{pmatrix}2&-1\\1&0\end{pmatrix}
\]
in coordinates $j,j+1$ and is the identity elsewhere.
It fixes $\mathbf 1=(1,\ldots,1)^t$ and preserves the
covector $\epsilon=(1,-1,\ldots,1)$.
Since $\epsilon(\mathbf1)=1$, quotienting by
$\mathbb Z\mathbf1$ is integrally isomorphic to the
invariant hyperplane $\ker\epsilon$.
Using quotient representatives with last coordinate zero,
the reduced matrix is
\[
        R_j=(B_j)_{1\ldots10,1\ldots10}
                    -\mathbf1_{10}(B_j)_{11,1\ldots10}.
                                                               \tag{B.62.4}
\]
It preserves the unimodular alternating matrix
\[
 J_{ij}=(-1)^{j-i+1}\quad(i<j),\qquad
 J_{ji}=-J_{ij},\qquad J_{ii}=0,
 \qquad \det J=1.                                           \tag{B.62.5}
\]
The program \nolinkurl{next_finite_symplectic/prepare.py}
checks these assertions over $\mathbb Z$, then the braid
relations and inverses modulo three.  It also checks
\[
       (R_1\cdots R_{10})^{11}=-I\pmod3.
\]
Projectivization kills the full twist and descends through
$PB_{11}/\langle\Delta^2\rangle\cong P$
as in (B.62.1).  Its pure image is all of
$\operatorname{PSp}_{10}(\mathbb F_3)$:
\citet[Theorem 5.1, with $g=5$, $b=1$, $m=6$]
{DE_Stylianakis2018} gives the surjection
$PB_{11}\to\operatorname{Sp}_{10}(\mathbb F_3)$,
and the full twist maps to $-I$.

\emph{The two actual subgroup images.}
Use the thirteen ordered tetrahedral meridians
$e_1,\ldots,e_{13}$ of B.56 and the six cyclic meridians
$f_0,\ldots,f_5$ of B.60, with exactly the common marked
base transport certified in B.61.
In particular the algebraic parameters are
$u_0=i(\varphi+\sqrt{\varphi+2})$ and $v_0=-\varphi^2$,
where $\varphi=(1+\sqrt5)/2$.
Both configurations normalize labels $2,4,6$ to
$\infty,0,1$, and the common finite order is
$(4,11,7,9,3,5,1,12,8,10,6)$.
The nineteen actual words, including the tetrahedral
$16$-crossing and cyclic $39$-crossing transports, are in
\nolinkurl{next_B59_image/common_base_braid_words.json}.
Their path and homotopy certificates were proved in B.61;
the present calculation does not conjugate either subgroup
independently after choosing this common base.

Put $A=q(H_T)$, $C=q(H_5)$ and $I=A\cap C$.
Exact finite multiplication gives
\[
 |A|=144,\qquad |C|=1440,\qquad I\cong C_3,
 \qquad |A'|=24,\qquad I\cap A'=1.                           \tag{B.62.6}
\]
Moreover $A/A'\cong C_6$.
Normalize its additive character
$\chi:A\to\mathbb Z/6$ by
$\chi(q(e_1e_2^{-1}))=1$.
The complete source values are
\[
 \bigl(\chi(q(e_i))\bigr)_{i=1}^{13}
       =(3,2,3,3,2,3,2,3,3,3,2,2,3),                       \tag{B.62.7}
\]
and $\chi(I)=\{0,2,4\}$.
In particular its restriction to $I$ is injective.

These are exhaustive matrix calculations, not order
estimates.  The program
\nolinkurl{next_finite_symplectic/enumerate.cpp}
enumerates both generated projective groups by closure
under their generators, representing every matrix over
$\mathbb F_3$ by the scalar multiple whose first nonzero
entry is one.
The separate program
\nolinkurl{next_finite_symplectic/analyse.py}
recomputes all $144^2$ products in $A$, its full derived
subgroup, the intersection, and the cyclic character.
The complete matrices, multiplication table and actual
standard output are retained.

Independently,
\nolinkurl{code/verify_B62_symplectic.py}
starts from the original nineteen braid words, works
instead on $\ker\epsilon$, and verifies (B.62.6), the
seven source relations below, and the integral change of
basis matching every discovery matrix.
Its passing output is
\nolinkurl{code/verify_B62_symplectic.out}.

\emph{All source-abelian directions and the central twist.}
Let $M$ be B.58's $55\times13$ integral matrix of all
pair linkings of the eleven normalized finite points.
It has rank six, and its integral kernel has the following
seven basis vectors, displayed as exponent-sum vectors
of actual free-group words:
\[
\begin{array}{lll}
 r_1=e_1e_2e_3e_4e_5e_6e_7,&
 r_2=e_3^{-1}e_8,&r_3=e_6^{-1}e_9,\\
 r_4=e_1^{-1}e_{10},&r_5=e_5^{-1}e_{11},&
 r_6=e_2^{-1}e_{12},\\
 r_7=e_4^{-1}e_{13}.&&
\end{array}                                                 \tag{B.62.8}
\]
If $B$ is their $13\times7$ column matrix, then $MB=0$
and the rows $7,\ldots,13$ of $B$ form the identity matrix.
The exact rank calculation therefore proves that these
columns form the full integral kernel, with no finite
index omitted.  Equation (B.62.7) gives
$\chi(q(r_j))=0$ for all seven $j$; direct finite-word
evaluation independently places all seven images in $A'$.

We must compare this matrix with the mapping-class
abelianization, rather than silently use disk braid
homology.  Capping gives
\[
 P^{\mathrm{ab}}
     =\mathbb Z^{55}/\mathbb Z\mathbf1_{55}.
                                                               \tag{B.62.9}
\]
Every chosen tetrahedral braid has zero linking on the
fixed pair $(4,6)$, whose coordinates are constantly
$0,1$.  Thus the corresponding row of $M$ is identically
zero.  This is checked in B.58's divisor certificate,
B.61's actual braid words, and both present matrix
audits.  If $Mv=\lambda\mathbf1_{55}$ over
$\widehat{\mathbb Z}$, that row forces $\lambda=0$.

Let $k\in\mathcal K$, and choose a preimage
$\hat a\in\widehat{H_T}$ under the canonical map onto
$\overline{H_T}\subseteq\widehat P$.
Such a preimage exists by compactness and density;
no comparison between a discrete test kernel and its
closure is needed here.
Write $v\in\widehat{\mathbb Z}^{13}$ for its source
abelianization.  B.58 says that $k$ has zero image in
$\widehat{P^{\mathrm{ab}}}$.
Equations (B.62.9) and the fixed row therefore give
$Mv=0$ in $\widehat{\mathbb Z}^{55}$.
Flatness of the torsion-free $\mathbb Z$-module
$\widehat{\mathbb Z}$ gives the exact equality
\[
 \ker\bigl(M:\widehat{\mathbb Z}^{13}
                    \to\widehat{\mathbb Z}^{55}\bigr)
                       =B\widehat{\mathbb Z}^{7}.            \tag{B.62.10}
\]
Equivalently, the displayed identity minor, together
with a nonzero rank-six minor of $M$, proves this directly
because multiplication by a nonzero integer is injective
on $\widehat{\mathbb Z}$.
The finite character $\chi\circ q$ factors through the
source abelianization and annihilates all seven columns
of $B$.  It thus kills $\hat a$.
Consequently $\widehat q(k)\in A'$.
But $k\in\overline{H_T}\cap\overline{H_5}$ also implies
$\widehat q(k)\in I$, so (B.62.6) proves (B.62.2).

\emph{Refinements and scope.}
Let $g\in P$ lie in the closure of (B.62.3).
B.52 gives $g=hj$ with $h\in H_T$, $j\in H_5$;
compactness gives a second factorization
$g=\hat h\hat j$ through the two completed refinements.
Then $h^{-1}\hat h=j\hat j^{-1}\in\mathcal K$.
By (B.62.2) this ambiguity lies in
$\ker\widehat q=\overline{N_3}\subseteq
\overline{L_1}\cap\overline{L_2}$.
Hence $h\in L_1$, $j\in L_2$, proving (B.62.3).

Nonnilpotence can already be checked without identifying
the whole symplectic image: a nilpotent group of order
$144=2^4\cdot3^2$ has an abelian Sylow-three subgroup,
so its derived subgroup has order a power of two.
The subgroup $A$ has derived order $24$, contradicting
that condition.  Thus this supplies refinements from
an actual nonnilpotent quotient, beyond the nilpotent
families of B.58--B.59.
The conclusion uses the entire source-abelian kernel in
(B.62.8); the mere finite intersection $I\ne1$ would not
have proved it.
\end{proof}
